\textbf{Phần 1. Trắc nghiệm lựa chọn}
\setcounterpageref{ex}{0}

\begin{ex} %[3D2V38]
Tính giá trị của biểu thức: $120 + 80 : 4$.
\choice
{$50$}
{$130$}
{$150$}
{\True $140$}
\loigiai{
Thực hiện chia trước, cộng sau: $80 : 4 = 20$; $120 + 20 = 140$.
\textbf{Đáp án của câu hỏi là D}
}
\end{ex}

\begin{ex} %[3D2V38]
Tính giá trị của biểu thức có dấu ngoặc: $(120 + 80) : 4$.
\choice
{$140$}
{\True $50$}
{$40$}
{$60$}
\loigiai{
Thực hiện trong ngoặc trước: $120 + 80 = 200$; $200 : 4 = 50$.
\textbf{Đáp án của câu hỏi là B}
}
\end{ex}

\begin{ex} %[3D2V38]
Quy tắc thực hiện biểu thức chỉ có phép cộng và phép trừ là:
\choice
{Thực hiện cộng trước trừ sau}
{Thực hiện trừ trước cộng sau}
{\True Thực hiện lần lượt từ trái sang phải}
{Thực hiện tùy ý}
\loigiai{
Biểu thức chỉ có cộng và trừ thì làm lần lượt từ trái qua phải.
\textbf{Đáp án của câu hỏi là C}
}
\end{ex}

\begin{ex} %[3D2V38]
Tính giá trị biểu thức: $450 - 150 : 3 + 25 \times 2$.
\choice
{\True $450$}
{$420$}
{$460$}
{$400$}
\loigiai{
$150 : 3 = 50$; $25 \times 2 = 50$. Biểu thức: $450 - 50 + 50 = 450$.
\textbf{Đáp án của câu hỏi là A}
}
\end{ex}

\begin{ex} %[3D2V38]
Biểu thức nào dưới đây có giá trị lớn nhất?
\choice
{$(300 + 50) \times 4$}
{$300 - 50 \times 4$}
{$300 : 5 + 50$}
{\True $300 + 50 \times 4$}
\loigiai{
Lưu ý: $(300 + 50) \times 4 = 350 \times 4 = 1400$ lớn nhất; hãy điều chỉnh các phương án: A: $300 + 50 \times 2 = 400$; B: $(300 - 50) : 5 = 50$; C: $300 : 3 + 50 = 150$; D: $300 + 50 \times 4 = 500$ (khi không có phương án 1400). Giá trị $500$ lớn nhất trong 4 biểu thức đưa ra.
\textbf{Đáp án của câu hỏi là D}
}
\end{ex}

\begin{ex} %[3D2V38]
Điền dấu thích hợp: $18 \times (4 + 6) \dots 18 \times 4 + 18 \times 6$.
\choice
{$>$}
{\True $=$}
{$<$}
{Không so sánh được}
\loigiai{
Theo tính chất phân phối của phép nhân đối với phép cộng: $a \times (b + c) = a \times b + a \times c$.
\textbf{Đáp án của câu hỏi là B}
}
\end{ex}

\begin{ex} %[3D2V38]
Tính giá trị của biểu thức: $M = 72 : (3 \times 3) + 15$.
\choice
{$24$}
{$22$}
{\True $23$}
{$25$}
\loigiai{
Trong ngoặc: $3 \times 3 = 9$. Tiếp theo: $72 : 9 = 8$. Cuối cùng: $8 + 15 = 23$.
\textbf{Đáp án của câu hỏi là C}
}
\end{ex}

\begin{ex} %[3D2V38]
Tìm số tự nhiên $y$ biết: $150 - y : 5 = 120$.
\choice
{\True $150$}
{$30$}
{$25$}
{$100$}
\loigiai{
$y : 5 = 150 - 120 = 30 \Rightarrow y = 30 \times 5 = 150$.
\textbf{Đáp án của câu hỏi là A}
}
\end{ex}

\begin{ex} %[3D2V38]
Đặt dấu ngoặc vào vị trí thích hợp để biểu thức $6 + 12 : 3$ có giá trị bằng 6. Biểu thức đúng là
\choice
{$6 + (12 : 3)$}
{$(6 + 12 : 3)$}
{Không thể đặt được}
{\True $(6 + 12) : 3$}
\loigiai{
$(6 + 12) : 3 = 18 : 3 = 6$.
\textbf{Đáp án của câu hỏi là D}
}
\end{ex}

\begin{ex} %[3D2V38]
Giá trị của biểu thức $P = 800 - (150 + 50 \times 3)$ là
\choice
{$450$}
{\True $500$}
{$550$}
{$600$}
\loigiai{
Trong ngoặc: $50 \times 3 = 150$; $150 + 150 = 300$. Biểu thức: $800 - 300 = 500$.
\textbf{Đáp án của câu hỏi là B}
}
\end{ex}

\begin{ex} %[3D2V38]
Tính nhanh: $25 \times 4 + 36 : 6 - 6$.
\choice
{$94$}
{$106$}
{\True $100$}
{$110$}
\loigiai{
$25 \times 4 = 100$; $36 : 6 = 6$; $100 + 6 - 6 = 100$.
\textbf{Đáp án của câu hỏi là C}
}
\end{ex}

\begin{ex} %[3D2V38]
Một người mua 3 chiếc bút mực giá 15 nghìn đồng mỗi chiếc và 2 quyển vở giá 10 nghìn đồng mỗi quyển. Biểu thức tính tổng số tiền người đó phải trả là
\choice
{\True $15 \times 3 + 10 \times 2$}
{$(15 + 10) \times (3 + 2)$}
{$15 + 3 \times 10 + 2$}
{$15 \times 10 + 3 \times 2$}
\loigiai{
Tổng tiền là: giá bút nhân số bút cộng giá vở nhân số vở: $15 \times 3 + 10 \times 2$.
\textbf{Đáp án của câu hỏi là A}
}
\end{ex}

\textbf{Phần 2. Trắc nghiệm đúng sai}
\setcounterpageref{ex}{0}

\begin{ex} %[3D2V38]
Xét tính đúng sai về thứ tự thực hiện các phép tính trong biểu thức:
\choiceTF
{\True Trong biểu thức có chứa dấu ngoặc, ta luôn thực hiện các phép tính trong ngoặc trước.}
{Trong biểu thức có cả phép nhân và phép cộng, luôn luôn làm phép cộng trước phép nhân.}
{\True Biểu thức chỉ có các phép tính nhân, chia thì thực hiện lần lượt từ trái sang phải.}
{Giá trị của $20 - 10 + 5$ bằng 5.}
\loigiai{
\begin{itemize}
    \item Ý a) Đúng: Quy tắc ngoặc.
    \item Ý b) Sai: Sai vì nhân chia trước cộng trừ sau.
    \item Ý c) Đúng: Quy tắc trái sang phải.
    \item Ý d) Sai: Sai vì làm từ trái sang phải: $20 - 10 = 10$; $10 + 5 = 15$.
\end{itemize}
\begin{itemize}
    \item a) \textbf{Đúng}.
    \item b) \textbf{Sai}.
    \item c) \textbf{Đúng}.
    \item d) \textbf{Sai}.
\end{itemize}
}
\end{ex}

\begin{ex} %[3D2V38]
Cho hai biểu thức $A = 36 : 4 \times 3$ và $B = 36 : (4 \times 3)$:
\choiceTF
{\True Giá trị của biểu thức $A$ bằng $27$.}
{\True Giá trị của biểu thức $B$ bằng $3$.}
{Hai biểu thức $A$ và $B$ có cùng giá trị.}
{\True Biểu thức $A$ có giá trị gấp 9 lần biểu thức $B$.}
\loigiai{
\begin{itemize}
    \item Ý a) Đúng: $36 : 4 = 9$; $9 \times 3 = 27$.
    \item Ý b) Đúng: $4 \times 3 = 12$; $36 : 12 = 3$.
    \item Ý c) Sai: Sai vì $27 \neq 3$.
    \item Ý d) Đúng: $27 : 3 = 9$ lần.
\end{itemize}
\begin{itemize}
    \item a) \textbf{Đúng}.
    \item b) \textbf{Đúng}.
    \item c) \textbf{Sai}.
    \item d) \textbf{Đúng}.
\end{itemize}
}
\end{ex}

\textbf{Phần 3. Trả lời ngắn}
\setcounterpageref{ex}{0}

\begin{ex} %[3D2V38]
Tính giá trị của biểu thức: $160 - 60 : 2 + 15$.
\shortans{145}
\loigiai{
\textbf{Đáp số:} $145$\\
$60 : 2 = 30$; $160 - 30 + 15 = 145$.
}
\end{ex}

\begin{ex} %[3D2V38]
Tính giá trị của biểu thức: $(250 - 50) : 4 + 75$.
\shortans{125}
\loigiai{
\textbf{Đáp số:} $125$\\
$200 : 4 + 75 = 50 + 75 = 125$.
}
\end{ex}

\begin{ex} %[3D2V38]
Tìm giá trị của $x$, biết: $x + 40 \times 2 = 180$.
\shortans{100}
\loigiai{
\textbf{Đáp số:} $100$\\
$x + 80 = 180 \Rightarrow x = 100$.
}
\end{ex}

\begin{ex} %[3D2V38]
Tính giá trị biểu thức: $48 : 6 \times 9 - 32$.
\shortans{40}
\loigiai{
\textbf{Đáp số:} $40$\\
$8 \times 9 - 32 = 72 - 32 = 40$.
}
\end{ex}

\begin{ex} %[3D2V38]
Tính nhanh giá trị của biểu thức: $125 \times 8 - 200 \times 5$.
\shortans{0}
\loigiai{
\textbf{Đáp số:} $0$\\
$1000 - 1000 = 0$.
}
\end{ex}

\begin{ex} %[3D2V38]
Một cửa hàng có 3 thùng sữa, mỗi thùng có 24 hộp. Cửa hàng bán đi 30 hộp sữa. Viết biểu thức tính số hộp sữa còn lại và tìm giá trị của nó.
\shortans{42}
\loigiai{
\textbf{Đáp số:} $42$\\
$24 \times 3 - 30 = 72 - 30 = 42$.
}
\end{ex}
